\chapter*{Conclusion Générale}
\addcontentsline{toc}{chapter}{Conclusion Générale}
\markboth{Conclusion Générale}{Conclusion Générale}

Ce projet a démontré, à travers une implémentation fonctionnelle et une validation
expérimentale rigoureuse, comment le Service Mesh Istio répond aux défis de sécurité
posés par les architectures microservices dans le contexte DevSecOps. La problématique
initiale portait sur l'incapacité du DAST traditionnel à inspecter et sécuriser le
trafic \textit{East-West} --- les communications interservices internes au cluster,
invisibles pour les outils de test conventionnels. Cette lacune a été adressée par
le déploiement d'une architecture en six couches de sécurité complémentaires : TLS
Gateway, RequestAuthentication JWT, AuthorizationPolicy RBAC, mTLS STRICT, SPIFFE
et NetworkPolicy.

Le prototype de système de paiement, composé de deux microservices déployés sur
Kubernetes/Minikube avec Istio, a servi de terrain d'expérimentation pour douze
scénarios d'attaque couvrant l'ensemble des six catégories du framework STRIDE.
Les résultats sont sans ambiguïté : onze scénarios sur douze ont été intégralement
bloqués par les mécanismes Istio, sans aucune modification du code applicatif.
Le douzième scénario (T10 --- flood sans rate limiting) a motivé l'implémentation
d'un EnvoyFilter de rate limiting (T11), illustrant la capacité d'extension
incrémentale du Service Mesh face à de nouvelles menaces.

L'apport le plus significatif de ce travail réside dans la démonstration que la
posture de sécurité est entièrement déterminée par les politiques appliquées au
mesh, indépendamment du code des applications. Le même binaire Go déployé sans
\texttt{mesh-config} est vulnérable ; avec \texttt{mesh-config}, il est protégé
contre onze catégories d'attaques distinctes. Cette séparation des préoccupations
--- sécurité portée par l'infrastructure, transparente pour le développeur ---
constitue l'essence même de l'apport DevSecOps du Service Mesh. Par ailleurs,
Kiali et Grafana rendent observable et auditable un trafic jusqu'alors invisible,
tandis que les logs Envoy garantissent un audit trail complet assurant la
non-répudiation de toutes les actions dans le mesh.

Ce travail présente néanmoins certaines limites. L'environnement Minikube, bien
que représentatif sur le plan fonctionnel, ne reproduit pas toutes les contraintes
d'un cluster de production multi-nœuds. La Gateway externe n'est pas configurée
en HTTPS dans ce prototype, une lacune à combler avant tout déploiement en
production. Plusieurs perspectives d'extension sont identifiées : l'intégration
des politiques Istio dans un pipeline CI/CD via OPA Gatekeeper, le couplage
d'OWASP ZAP au mesh pour automatiser les scans internes, l'ajout d'un moteur
de détection d'anomalies comportementales tel que Falco, et l'extension vers une
architecture Zero Trust multi-cluster.

Au-delà des aspects techniques, ce projet illustre une transformation culturelle
fondamentale du DevSecOps : la sécurité n'est plus un obstacle ou une étape
finale du cycle de développement, mais une propriété émergente de l'infrastructure.
Lorsque les politiques de sécurité sont déclaratives, versionnables et appliquées
automatiquement par le mesh, elles deviennent aussi naturelles qu'un déploiement
Kubernetes ordinaire. C'est cette philosophie --- \textit{sécurité by default},
\textit{sécurité as code}, \textit{sécurité sans friction} --- qui constitue la
promesse centrale du Service Mesh dans l'écosystème DevSecOps.